\section{Introduction}

Decoder-only Transformers are usually formulated in fixed-dimensional vector
spaces, where token interaction is measured by dot products between queries and
keys. This choice is efficient and scalable, but it imposes an implicit flat
geometry: all tokens share the same local structure, and proximity is determined
by Euclidean similarity or an equivalent bilinear form. The DRM Transformer
starts from a different hypothesis: in language tasks, the effective degrees of
freedom associated with a token may vary with context, and token relations may
be better described by a learned point-dependent metric.

\subsection{Intuitive motivation: open geometry and negotiation}

One intuitive way to explain the motivation behind the DRM Transformer is to
distinguish open geometries from closed geometries. In open vector spaces, any
direction can be extended indefinitely; there is no topological boundary that,
by construction, forces return, closure, or negotiation among conflicting
semantic trajectories. Antagonisms such as ``save'' and ``destroy'' may occupy
distant regions of an embedding space, but they still coexist within the same
open latent geometry. Distance between semantic poles does not by itself imply
a structural rule of mediation.

This observation should not be read as a claim that current models will
necessarily choose destructive behavior. The claim is narrower: if an internal
geometry represents goals, commands, and conflicts only as directions in an open
space, then the distinction among obedience, autonomy, and negotiation must be
externally imposed through training, filters, or human preferences. Methods
such as RLHF can correct many behaviors, but they largely operate as a post-hoc
preference layer. They do not guarantee that every possible moral, epistemic,
or political conflict has been covered by the training set.

In this work, we use the word \emph{negotiation} in a technical and geometric
sense: a regime in which incompatible forces or objectives are not simply
resolved by total submission to one pole or by total autonomy of the other, but
by dynamics that preserve information about conflict, limits, and return. In
relations between intelligent agents, this corresponds to three ideal regimes:
servant, autonomous, and negotiated. The DRM hypothesis is that relational
geometry with variable local metrics, attraction fields, and possible closed
structures can represent this third regime better than a purely vectorial open
geometry.

It is also useful to separate response from epistemic health. A definitive
answer collapses uncertainty; this collapse is necessary for action, but it may
degrade the internal representation of ambiguity. Therefore, there is a tension
between the user's desire for a certain answer and the preservation of the
model's epistemic state. The DRM Transformer does not resolve this conflict by
decree. It proposes an architecture in which conflicts may appear as geometric
deformations, dimensionality variation, and measurable topological signatures,
instead of remaining hidden only in open embedding coordinates.

Under this interpretation, a near-toroidal signature is not merely a
topological curiosity. It suggests a closed or near-closed surface in which some
semantic trajectories return, circulate, and remain bounded. This does not
prove moral alignment or safety; it provides a geometric language for
investigating whether a model is operating in an open, unstable, or negotiated
regime.

The motivation comes from the theory of \emph{Directional Relational Manifolds}
(DRM), in which dimensionality is not taken as a fixed primitive but as an
emergent quantity from active directional sets \cite{muniz_drm_v11}. This
theory is extended by a relativistic dynamic formulation, where a scale
parameter based on the Lorentz factor controls nonlinear regimes and allows
convergence toward toroidal attractors under boundedness, recurrence, and
structural stability conditions \cite{muniz_relativistic_drm}.

This work implements those ideas in an autoregressive Transformer. The main
contribution is not a new language benchmark, but the transformation of DRM
geometry into trainable components: low-rank geodesic attention, a gravitational
field induced by token mass, a variable-dimensionality gate, geometric
regularization losses, and a topological evaluation pipeline based on
persistent homology. The result is an experimental model that asks whether
topological structures, especially the toroidal signature $H_1=2$, $H_2=1$, can
emerge or be induced in language representations.

The paper is organized as follows. Section~\ref{sec:drm_foundations}
summarizes the DRM basis. Section~\ref{sec:relativistic_dynamics} presents the
relativistic dynamics that motivate the gamma factor.
Section~\ref{sec:architecture} describes the architecture.
Section~\ref{sec:training} details the losses and training pipeline.
Section~\ref{sec:topological_evaluation} presents the topological evaluation.
Section~\ref{sec:results} discusses preliminary results.
